\section{Conclusion}
\label{sec:conclusion}

This work investigated error-based active disturbance rejection control for wrist tremor suppression in Parkinson's disease while accounting for upper limb dynamics. A 3-DoF mass-spring-damper model enabled the evaluation of control performance under disturbances arising from shoulder and elbow tremors. The results show that EADRC strategies, including EBMFLC and ZPLP variations, achieve high tremor power suppression with lower control effort than the more demanding DE-PID approach. 
%
\reviewchange{}{Of the two, EADRC+ZPLP achieved higher tremor suppression with smoother control signals.}
%
In particular, EADRC+ZPLP also reached an average tremor power suppression rate of 56.97\% in dynamic scenarios, outperforming classical controllers.

EADRC demonstrated strong robustness with both voluntary motion estimators used, maintaining consistent performance under variations of up to 50\% in joint stiffness and damping parameters. This indicates its suitability for practical applications where biomechanical properties are uncertain or time-varying. 
%
EADRC's ease of implementation and flexibility not withstanding, voluntary motion estimation is still key in avoiding the inclusion of lower-frequency tremors in the total estimated disturbance.
%
Future work will consider the inclusion of more control strategies from the literature, 
as well as the investigation of alternative voluntary motion estimation methods, and use of different upper limb models.